%! TeX program = lualatex
% Copyright (c) 1996, 2004, 2006, 2009, 2014, 2016, 2019, 2023
% Tama Communications Corporation. All rights reserved.
%
% Redistribution and use in source and binary forms, with or without
% modification, are permitted provided that the following conditions
% are met:
% 1. Redistributions of source code must retain the above copyright
%    notice, this list of conditions and the following disclaimer.
% 2. Redistributions in binary form must reproduce the above copyright
%    notice, this list of conditions and the following disclaimer in the
%    documentation and/or other materials provided with the distribution.
%
% THIS SOFTWARE IS PROVIDED BY THE AUTHOR AND CONTRIBUTORS ``AS IS'' AND
% ANY EXPRESS OR IMPLIED WARRANTIES, INCLUDING, BUT NOT LIMITED TO, THE
% IMPLIED WARRANTIES OF MERCHANTABILITY AND FITNESS FOR A PARTICULAR PURPOSE
% ARE DISCLAIMED.  IN NO EVENT SHALL THE AUTHOR OR CONTRIBUTORS BE LIABLE
% FOR ANY DIRECT, INDIRECT, INCIDENTAL, SPECIAL, EXEMPLARY, OR CONSEQUENTIAL
% DAMAGES (INCLUDING, BUT NOT LIMITED TO, PROCUREMENT OF SUBSTITUTE GOODS
% OR SERVICES; LOSS OF USE, DATA, OR PROFITS; OR BUSINESS INTERRUPTION)
% HOWEVER CAUSED AND ON ANY THEORY OF LIABILITY, WHETHER IN CONTRACT, STRICT
% LIABILITY, OR TORT (INCLUDING NEGLIGENCE OR OTHERWISE) ARISING IN ANY WAY
% OUT OF THE USE OF THIS SOFTWARE, EVEN IF ADVISED OF THE POSSIBILITY OF
% SUCH DAMAGE.
%
%	rireki.tex	Version 3.1
%
%	URL: https://www.tamacom.com/rireki-j.html
%
%----------------------------------------------------------------------------
%
% 多摩通信社からのお知らせ
%
% 「履歴書猿人」(https://www.tamacom.com/engine/rireki2/) という履歴書作成のための
% Webアプリケーションを公開しております。TeX の知識なしにTeX の履歴書が作成でき、
% 様々なカスタマイズも可能です。こちらもご利用いただけると幸いです。

%----------------------------------------------------------------------------
%
% ヘッダー
%
\documentclass[b5j]{ltjsarticle}
\usepackage[deluxe,nfssonly]{luatexja-preset}
\usepackage{graphicx}
\usepackage{../../../../Styles/Rireki-Styles/Rireki-Style-JP}

\usepackage{datetime2}
\usepackage[
pdftitle={履歴書},
pdfauthor={アンドレ・モッシ},
pdfcreator={LaTeXとNeoVim}
]{hyperref}

% オプション
%
% 下記のオプションが利用可能です。

\空行挿入		% 学歴と職歴の間に空行を挿入します
%\性別欄なし		% 性別欄を削除します
%\写真欄なし		% 写真欄を削除します
\begin{document}

% ID情報

\姓{\LARGE Mossi Milla}
\名{\LARGE Roberto Andre}
\姓読み{もっし　みら}
\名読み{ろばーと　あのどれ}
\性別{男}					% 男|女
\誕生日{2002年5月27日}
\現在日付{\the\year 年\the\month 月\the\day 日}
\年齢{\number\numexpr\the\year - 2002\relax}

% 顔写真
%
% 画像ファイルにはEPS フォーマット・縦横比4:3 のものをご使用ください。
% 縦を4cm に調整し、縦横比を変更せずに印刷します。
% 次のように指定します。
% \顔写真{photo.jpg}

\顔写真{../../../../Assets/2025PFPic.png}

% 現住所

\現住所郵便番号{16541}
\現住所{アメリカ合衆国 ペンシルベニア州 エリー市 109 University Square}
\現住所読み{ゆにばーしてぃ　すくえあ　えりー　ぺんしるばにあ}
\現住所市外局番{}
\現住所電話番号{}
\現住所呼び出し{}

% 連絡先

\連絡先郵便番号{}
\連絡先{\tt mossiroberto0392@gmail.com}
\連絡先読み{}
\連絡先市外局番{}
\連絡先電話番号{+1 814-790-1591}
\連絡先呼び出し{}

% 学歴、職歴
%
% 学歴、職歴を年月順に列挙してください。合計20個まで記入出来ます。
% 20個を超える部分は印刷されませんので、ご注意ください。
% 印刷順は、学歴=>職歴の順になります。

\学歴{2015}{4}{Dowal School　入学}      % {年}{月}{内容}
\学歴{2020}{3}{Dowal School　卒業}
\学歴{2021}{8}{Gannon University コンピュータサイエンス学科　入学}
\学歴{2025}{5}{Gannon University コンピュータサイエンス学科　卒業見込み}
\学歴{2019}{4}{コンピュータクラブ会長 就任}
\学歴{2020}{3}{技術コンペティション優勝（ゲーム開発）}
\学歴{2025}{10}{ISI日本語学校　入学}

\職歴{2023}{6}{Dinant Company インターンとして入社}            % {取得年}{取得月}{資格}
\職歴{2023}{7}{Dinant Company インターン期間満了退社}

% 資格
%
% 資格を取得年月順に列挙してください。9つまで記入できます。
% 9つを超える部分は印刷されませんので、ご注意ください。

\資格{2023}{8}{ホンジュラスの運転免許証}            % {取得年}{取得月}{資格}
\資格{2022}{12}{ハーバード大学 CS50's Introduction to Computer Science終了}
\資格{2022}{5}{エクストリームネットワークス Introduction to Future Networks終了}

% 個人情報
%
% 志望の動機と本人希望記入欄はlatex のコマンドを記述できます。

\志望の動機{
14歳の頃にゲーム開発をきっかけとしてプログラミングに興味を持ち、以来C++およびUnreal Engineを
中心に継続的に技術を磨いてまいりました。大学ではコンピュータサイエンスを専攻し、強化学習および
神経進化を用いたマルチエージェントシミュレーションの研究に取り組みました。

研究では、Unreal Engine上にシミュレーション環境を構築し、エージェントの学習過程を世代単位で分析することで、
行動パターンの進化を評価しました。本研究はASEEおよびSigma Xiにて受理され、技術的な専門性と継続的な探究心を
評価していただきました。

また、インターンシップではWebアプリケーション開発に従事し、実際のビジネス環境におけるチーム開発および
要件定義の重要性を学びました。個人開発においても、Neovim向けプラグインの開発などを通じて、設計力と
継続的改善を意識した開発を行っております。

これまで培ったC++およびシミュレーション技術を活かし、実社会に価値を提供できるシステム開発に貢献したいと
考えております。常に学習を続けながら、チームの一員として長期的に成長していきたいと存じます。
}
\本人希望記入欄{
\begin{itemize}
	\item \textbf{ポートフォリオウェブサイト} https://andremossi.vercel.app/ \\
	      私のストーリーや、私が開発したプロジェクトや、キャリアに関連する仕事の情報を含むウェブサイトです。

	\item \textbf{リンクツリーウェブサイト} https://andremossi-linktree.vercel.app/ \\
	      私のソーシャルメディアアカウントへのリンクが含まれているリンクツリーウェブサイトです。

	\item \textbf{Github} https://github.com/AndreM222 \\
	      私のプロジェクトのコードを公開しています。

\end{itemize}
}

\サイン{Roberto Andre Mossi Milla}

\end{document}
